% ---------------------------------------------------------------------------
% Training-free inpainting with a diffusion prior: implementation note
% Designed to be \input{} into the paper; requires amsmath, booktabs.
% Bibliography keys are provided in inpainting_methods.bib.
% ---------------------------------------------------------------------------

\section{Training-free inpainting with a diffusion prior}
\label{sec:inpainting}

\subsection{Diffusion prior and notation}
\label{sec:inpainting:notation}

All inpainting algorithms considered in this work operate on a single
pre-trained denoising diffusion probabilistic model
(DDPM)~\cite{ho2020ddpm} per centrality class, used \emph{as is}: no
component is retrained or fine-tuned for the inpainting task.  The model
acts in the normalized (logarithmic) representation of the calorimeter
image, $x = \ln\!\bigl(\max(E_T, E_{\min})\bigr)$ with
$E_{\min} = 10^{-3}$~GeV, on images $x \in \mathbb{R}^{1\times 24\times 64}$
in the $(\eta,\phi)$ plane.

The forward process is the standard variance-preserving Markov chain with
$T$ steps,
\begin{equation}
  q(x_t \mid x_{t-1}) = \mathcal{N}\!\bigl(x_t;\,
      \sqrt{1-\beta_t}\, x_{t-1},\; \beta_t \mathbf{I}\bigr),
  \qquad
  q(x_t \mid x_0) = \mathcal{N}\!\bigl(x_t;\,
      \sqrt{\bar\alpha_t}\, x_0,\; (1-\bar\alpha_t)\,\mathbf{I}\bigr),
  \label{eq:forward}
\end{equation}
with $\alpha_t = 1-\beta_t$ and $\bar\alpha_t = \prod_{i\le t}\alpha_i$.
The schedule replicates the training configuration exactly: $T = 8000$ and
a linear schedule $\beta_t$ interpolating from $\beta_1 = b/T$ to
$\beta_T = 1000\,b/T$, with $b = 0.02$ ($0.10$) for the 0--10\%
(40--50\%) centrality model.  The network is an
$\epsilon$-prediction U-Net, $\hat\epsilon = \epsilon_\theta(x_t, t)$, from
which the denoised estimate follows by Tweedie's formula,
\begin{equation}
  \hat{x}_0(x_t)
  = \frac{x_t - \sqrt{1-\bar\alpha_t}\;\epsilon_\theta(x_t, t)}
         {\sqrt{\bar\alpha_t}} .
  \label{eq:tweedie}
\end{equation}

For sampling efficiency the reverse process is run on a subsampled time
grid $0 = \tau_0 < \tau_1 = 1 < \dots < \tau_S = T$ obtained by rounding a
linear spacing of $S$ points ($S = 1000$ unless stated otherwise); the
network is always conditioned on the \emph{original} time indices
$\tau_k$.  Writing $t = \tau_k$ and $s = \tau_{k-1}$ for one reverse step,
the effective per-step noise level is
$\beta_{(k)} \equiv 1 - \bar\alpha_t/\bar\alpha_s$, and the DDPM posterior
used by the ancestral sampler is
\begin{equation}
  q(x_s \mid x_t, x_0)
  = \mathcal{N}\!\left(x_s;\;
      \frac{\sqrt{\bar\alpha_s}\,\beta_{(k)}}{1-\bar\alpha_t}\, x_0
      + \frac{\sqrt{\bar\alpha_t/\bar\alpha_s}\,(1-\bar\alpha_s)}
             {1-\bar\alpha_t}\, x_t,\;
      \tilde\beta_{(k)} \mathbf{I}\right),
  \qquad
  \tilde\beta_{(k)} = \frac{(1-\bar\alpha_s)\,\beta_{(k)}}{1-\bar\alpha_t}.
  \label{eq:posterior}
\end{equation}
Because $\bar\alpha_{\tau_0} = \bar\alpha_0 = 1$, the final reverse step
collapses, $q(x_0 \mid x_{\tau_1}, \hat x_0) = \delta(x_0 - \hat x_0)$;
this property is used below to guarantee exact data consistency of the
final sample.  Unconditional generation initializes
$x_T \sim \mathcal{N}(0, (1-\bar\alpha_T)\mathbf{I})$ and iterates
Eqs.~(\ref{eq:tweedie})--(\ref{eq:posterior}) with $x_0 \to \hat x_0$.

\subsection{Noise-free inpainting setting}
\label{sec:inpainting:setting}

A dead detector region is described by a binary mask
$m \in \{0,1\}^{24\times 64}$, with $m_i = 1$ on live (observed) towers
and $m_i = 0$ on dead towers.  The measurement is noise-free,
\begin{equation}
  y = m \odot x_0^{\mathrm{true}},
  \qquad \sigma_y = 0,
  \label{eq:measurement}
\end{equation}
i.e.\ the degradation operator is $A = \operatorname{diag}(m)$ and the
inference target is the posterior
$p\bigl(x_0 \mid m \odot x_0 = y\bigr)$ under the DDPM prior.  Two
simplifications specific to this setting are used throughout.  First,
$A$ is its own Moore--Penrose pseudoinverse on its range,
$A^\dagger = A^{\!\top} = A$, so no singular-value machinery is required:
spectral space coincides with pixel space, with unit singular values on
observed and zero on dead pixels.  Second, all branches of the algorithms
associated with measurement noise ($\sigma_y > 0$) drop out identically.
In this noiseless-masking limit the five algorithms below---although
derived from different principles (heuristic projection, null-space
decomposition, spectral variational inference, and two flavors of
guidance by the measurement-consistency gradient)---approximate the
\emph{same} posterior, which makes their head-to-head calibration
comparison meaningful.

Posterior sampling is performed by batching $n$ independent reverse
trajectories per image ($n = 50$ in the main study); each algorithm
consumes one measurement pair $(y, m)$ and returns $n$ samples of the
full image, of which only the dead region is retained for analysis.

\subsection{Algorithms}
\label{sec:inpainting:algorithms}

Throughout this section one reverse step from $t = \tau_k$ to
$s = \tau_{k-1}$ is described; $\xi, \xi', \dots \sim \mathcal{N}(0,
\mathbf{I})$ denote independent standard normal draws, and
$\hat x_0 \equiv \hat x_0(x_t)$, $\hat\epsilon \equiv
\epsilon_\theta(x_t, \tau_k)$ from Eq.~(\ref{eq:tweedie}).

\subsubsection{RePaint}
\label{sec:inpainting:repaint}

RePaint~\cite{lugmayr2022repaint} alternates an unconditional reverse
step on the dead region with an independent forward-diffusion of the
known content to the \emph{matching} noise level:
\begin{align}
  x_s^{\mathrm{known}} &\sim
      q(x_s \mid x_0 = y)
      = \mathcal{N}\!\bigl(\sqrt{\bar\alpha_s}\, y,\;
                           (1-\bar\alpha_s)\,\mathbf{I}\bigr),
  \label{eq:repaint-known}\\
  x_s^{\mathrm{dead}} &\sim
      q(x_s \mid x_t, \hat x_0)
      \quad\text{[Eq.~(\ref{eq:posterior})]},
  \label{eq:repaint-unknown}\\
  x_s &= m \odot x_s^{\mathrm{known}}
        + (1-m) \odot x_s^{\mathrm{dead}} .
  \label{eq:repaint-combine}
\end{align}
We emphasize the time index in Eq.~(\ref{eq:repaint-known}): both branches
must live at level $s = \tau_{k-1}$ \emph{after} the step.  Noising the
known region to level $\tau_k$ instead---a plausible-looking off-by-one---
leaves it systematically one step too noisy at every iteration and biases
the reconstruction at the dead-region boundary; we verified that this
variant degrades calibration and use the correct $\tau_{k-1}$ convention.

To harmonize the two branches, RePaint applies \emph{resampling}: steps
(\ref{eq:repaint-known})--(\ref{eq:repaint-combine}) are repeated $U$
times ($U = 10$ by default), jumping back one step between repetitions
with the forward kernel
$x_t \sim q(x_t \mid x_s) = \mathcal{N}(\sqrt{\bar\alpha_t/\bar\alpha_s}\,
x_s,\, \beta_{(k)}\mathbf{I})$.\footnote{Algorithm~1 of
Ref.~\cite{lugmayr2022repaint} writes the jump-back kernel with
$\beta_{t-1}$; the kernel consistent with the forward process
(\ref{eq:forward}) uses the transition variance of the step being
re-traversed, which is what we implement.}
At the final step $\bar\alpha_{\tau_0} = 1$, so
Eq.~(\ref{eq:repaint-known}) returns $y$ exactly: the known region of the
output equals the measurement identically.

\subsubsection{DDNM}
\label{sec:inpainting:ddnm}

The Denoising Diffusion Null-space Model~\cite{wang2022ddnm} rectifies
the range-space component of the denoised estimate at every step.  Any
solution consistent with $y = A x_0$ can be written
$x_0 = A^\dagger y + (\mathbf{I} - A^\dagger A)\, \bar x$; for noiseless
masking the projection of $\hat x_0$ onto this affine subspace is simply
\begin{equation}
  \hat x_0^{\,\mathrm{proj}} = m \odot y + (1-m) \odot \hat x_0 ,
  \label{eq:ddnm-proj}
\end{equation}
after which the standard ancestral step, Eq.~(\ref{eq:posterior}) with
$x_0 \to \hat x_0^{\,\mathrm{proj}}$, is taken.  No noisy-case
corrections (the $\Sigma$-scaling of DDNM$^{+}$) are needed for
$\sigma_y = 0$.  By the collapse property of the final step, the output
equals $\hat x_0^{\,\mathrm{proj}}$ at $k=1$, so data consistency of the
known region is again exact.

\subsubsection{DDRM}
\label{sec:inpainting:ddrm}

Denoising Diffusion Restoration Models~\cite{kawar2022ddrm} define a
variational posterior in the spectral domain of $A$.  For masking the
spectral basis is the pixel basis, and with $\sigma_y = 0$ the general
update [Eqs.~(7)--(8) of Ref.~\cite{kawar2022ddrm}] reduces, in the
variance-preserving variables used here, to
\begin{align}
  \text{dead pixels:}\quad
  x_s &= \sqrt{\bar\alpha_s}\,\hat x_0
        + \sqrt{(1-\eta^2)(1-\bar\alpha_s)}\;\hat\epsilon
        + \eta \sqrt{1-\bar\alpha_s}\;\xi ,
  \label{eq:ddrm-dead}\\
  \text{observed pixels:}\quad
  x_s &= \sqrt{\bar\alpha_s}\,\bigl[(1-\eta_b)\,\hat x_0
        + \eta_b\, y\bigr]
        + \sqrt{1-\bar\alpha_s}\;\xi' ,
  \label{eq:ddrm-obs}
\end{align}
with hyperparameters $\eta = 0.85$ and $\eta_b = 1$ at their published
defaults.  Initialization follows the same pattern at $t = T$:
observed pixels start from
$\mathcal{N}(\sqrt{\bar\alpha_T}\, y, (1-\bar\alpha_T)\mathbf{I})$ and
dead pixels from $\mathcal{N}(0, (1-\bar\alpha_T)\mathbf{I})$.
Note that Eq.~(\ref{eq:ddrm-dead}) is DDIM-like but references the
\emph{marginal} noise scale $\sqrt{1-\bar\alpha_s}$ rather than the
ancestral posterior width $\smash{\tilde\beta_{(k)}^{1/2}}$---a
DDRM-specific choice.  With $\eta_b = 1$ the observed pixels are
independently re-noised copies of $y$ at every level, and at the final
step ($\bar\alpha_0 = 1$, vanishing noise) the output known region equals
$y$ exactly while dead pixels return $\hat x_0$.

\subsubsection{MCG}
\label{sec:inpainting:mcg}

Manifold-constrained gradients~\cite{chung2022mcg} augment the
RePaint-style projection with a measurement-consistency gradient taken
\emph{through the denoiser network}.  With
$\mathcal{R}(x_t) = \bigl\lVert m \odot \bigl(y - \hat x_0(x_t)\bigr)
\bigr\rVert_2$ (the $\ell_2$ norm, not its square, following the official
implementation---this makes the correction self-normalizing), one step
reads
\begin{align}
  g &= \nabla_{x_t}\, \mathcal{R}(x_t),\\
  x_s^{\mathrm{dead}} &= x_s' - \alpha\, g,
  \qquad x_s' \sim q(x_s \mid x_t, \hat x_0),\\
  x_s &= m \odot x_s^{\mathrm{known}} + (1-m) \odot x_s^{\mathrm{dead}},
\end{align}
with $x_s^{\mathrm{known}}$ as in Eq.~(\ref{eq:repaint-known}) (same
$\tau_{k-1}$ convention) and step size $\alpha = 1$.  Note the division
of roles: the residual $\mathcal{R}$ is evaluated on the \emph{observed}
pixels, while the correction acts on the \emph{dead} pixels---the two
regions being coupled through the network Jacobian, which is what allows
observed-data consistency to steer the reconstruction inside the dead
region (the known pixels themselves are overwritten by the projection).
The gradient requires one backward pass through $\epsilon_\theta$ per
step; the prior network weights remain frozen, and differentiation is
with respect to $x_t$ only.  Exact data consistency of the known region
follows from the projection as for RePaint.

\subsubsection{$\Pi$GDM}
\label{sec:inpainting:pigdm}

Pseudoinverse-guided diffusion models~\cite{song2023pigdm} approximate
the intractable guidance term $\nabla_{x_t} \log p_t(y \mid x_t)$ using a
Gaussian surrogate for $p(x_0 \mid x_t)$.  For a noiseless measurement the
guidance reduces to a vector--Jacobian product with the mask-projected
residual [Eq.~(8) of Ref.~\cite{song2023pigdm}; the $r_t^2$ prefactor
cancels exactly in the update for $\sigma_y = 0$],
\begin{equation}
  g = \left(\frac{\partial \hat x_0(x_t)}{\partial x_t}\right)^{\!\!\top}
      \bigl[\, m \odot (y - \hat x_0)\,\bigr],
  \label{eq:pigdm-vjp}
\end{equation}
evaluated with one backward pass through the network (the residual is
treated as a constant cotangent).  The update is a DDIM($\eta$)
step~\cite{song2021ddim} plus the guidance scaled by
$\sqrt{\bar\alpha_t}$, as required for variance-preserving
parametrizations [Algorithm~1 of Ref.~\cite{song2023pigdm}]:
\begin{equation}
  x_s = \sqrt{\bar\alpha_s}\, \hat x_0
        + c_2\, \hat\epsilon
        + c_1\, \xi
        + \sqrt{\bar\alpha_t}\; g,
  \qquad
  c_1 = \eta\,\tilde\beta_{(k)}^{1/2},
  \quad
  c_2 = \sqrt{1 - \bar\alpha_s - c_1^2},
  \label{eq:pigdm-update}
\end{equation}
with $\eta = 1$ (for which the first three terms are equal in law to the
ancestral step).  Unlike the other four algorithms, $\Pi$GDM enforces
consistency during sampling only through the guidance term, so the known
region of its raw output matches $y$ approximately rather than
identically.  For the noise-free setting we project the known region to
$y$ at the final step; since this occurs after the last network
evaluation, it is pure post-processing---it cannot affect the sampling
dynamics or the dead-region pixels on which the analysis is based.

\subsection{Implementation and validation}
\label{sec:inpainting:implementation}

\begin{table}[t]
  \centering
  \caption{Summary of the implemented inpainting algorithms, specialized
  to noise-free masking.  ``Backprop'' indicates a backward pass through
  the frozen $\epsilon$-network per reverse step; ``exact~$y$'' indicates
  that the known region of the final sample equals the measurement
  identically.  NFEs are network function evaluations per posterior
  sample for $S$ reverse steps.
  ($^{\dagger}$: via a final-step projection applied after the last
  network evaluation; the published algorithm is only approximately
  consistent.)}
  \label{tab:inpainters}
  \begin{tabular}{lllll}
    \toprule
    Algorithm & Hyperparameters & Backprop & Exact $y$ & NFEs \\
    \midrule
    RePaint~\cite{lugmayr2022repaint} & $U = 10$ & no  & yes & $U(S{-}1){+}1$ \\
    DDNM~\cite{wang2022ddnm}          & ---      & no  & yes & $S$ \\
    DDRM~\cite{kawar2022ddrm}         & $\eta = 0.85$, $\eta_b = 1$ & no & yes & $S$ \\
    MCG~\cite{chung2022mcg}           & $\alpha = 1$ & yes & yes & $S$ (+$S$ bwd) \\
    $\Pi$GDM~\cite{song2023pigdm}     & $\eta = 1$   & yes & yes$^{\dagger}$ & $S$ (+$S$ bwd) \\
    \bottomrule
  \end{tabular}
\end{table}

All five algorithms share the sampler infrastructure: the exact training
schedule reconstructed from the model configuration, the subsampled time
grid of Sec.~\ref{sec:inpainting:notation}, evaluation of
$\epsilon_\theta$ under \texttt{bfloat16} autocast with all schedule and
update arithmetic in \texttt{float32}, and a dedicated
counter-based random stream per truth image (making every study run
bitwise reproducible and restartable).  Truth images are drawn from the
same DDPM used as the prior, so the ground truth is, by construction, a
sample from the prior---the setting in which posterior-calibration
diagnostics cleanly isolate algorithmic approximation error.
Table~\ref{tab:inpainters} summarizes the algorithm-specific choices.

The implementation is validated in two independent ways.  First, the
schedule, subsampling, and posterior coefficients are verified
numerically against the training code of Ref.~\cite{torbunov2024calo}
to single precision.  Second, each algorithm is run on an analytically
tractable problem: for an i.i.d.\ Gaussian prior
$x_0 \sim \mathcal{N}(\mu_0, \sigma_0^2 \mathbf{I})$ the optimal
$\epsilon$-predictor is known in closed form and the true posterior on
dead pixels is again $\mathcal{N}(\mu_0, \sigma_0^2)$, independent of the
observed pixels.  RePaint, DDNM, MCG, and $\Pi$GDM reproduce this
posterior to within sampling accuracy, and the algorithms with hard
projection preserve the known region exactly (to floating-point
precision).  DDRM at its default $\eta = 0.85$ reproduces the posterior
mean but \emph{under-disperses} the posterior (standard-deviation ratio
$\approx 0.86$ in this test, approaching $1$ as $\eta \to 0$ and
$\approx 0.70$ at $\eta = 1$, independent of $S$)---an intrinsic property
of its variational posterior in the noise-free limit rather than an
implementation artifact, and one that the calibration study of
Sec.~\ref{sec:results} is designed to expose.

\section{Statistical validation methodology}
\label{sec:validation}

\subsection{Evaluation frame}
\label{sec:validation:frame}

Truth images are drawn from the same DDPM that serves as the inpainting
prior.  The ground truth is therefore, by construction, a sample from the
prior, and for a noiseless measurement the pair
$(x_0^{\mathrm{true}},\, y = m \odot x_0^{\mathrm{true}})$ realizes
exactly the self-consistency setting of simulation-based
calibration~\cite{talts2018sbc}: if an algorithm sampled the posterior
$p(x_0 \mid y)$ exactly, truth and posterior samples would be
exchangeable, and any measurable departure from the calibration
diagnostics below is attributable purely to algorithmic approximation
error (there is no model misspecification by construction).

For each truth image $i = 1, \dots, N$ the algorithm under test draws
$J$ i.i.d.\ posterior samples $\{x^{(i,j)}\}_{j=1}^{J}$ ($J = 50$,
$N = 1000$ in the main study); all diagnostics are evaluated on the
dead-region pixels only.  Perceptual or point-estimate image metrics
(FID, LPIPS, PSNR, SSIM, Pearson correlation of the posterior mean) are
deliberately not used: the scientific target is the posterior
\emph{distribution} over physical quantities, not the fidelity of a
single reconstruction.  In particular, the posterior mean is an
intentionally variance-suppressed point estimate; regressing it against
truth necessarily yields a slope below unity.  This shrinkage slope is
reported as a diagnostic of the posterior width but is not scored.

\subsection{Per-pixel rank statistics (SBC)}
\label{sec:validation:sbc}

For every image $i$ and dead pixel $p$, the rank of truth among the
posterior samples is
\begin{equation}
  r_{ip} \;=\; \#\bigl\{\, j : x^{(i,j)}_p < x^{\mathrm{true}}_{ip} \bigr\}
  \;+\; u_{ip},
  \qquad
  u_{ip} \sim \mathcal{U}\bigl\{0, \dots,
      \#\{ j : x^{(i,j)}_p = x^{\mathrm{true}}_{ip}\}\bigr\},
  \label{eq:sbc-rank}
\end{equation}
where the randomized tie-breaking term $u_{ip}$ handles exact ties
(which occur at the $E_{\min}$ clipping floor).  Under exact calibration
$r_{ip}$ is uniform on $\{0, \dots, J\}$~\cite{talts2018sbc}.  Ranks are
invariant under strictly monotone transformations, so the GeV and
log-space analyses coincide.  We report the pooled rank histogram with a
$\chi^2$ uniformity $p$-value and, separately, per-pixel $p$-value maps
across the dead-region box.  The histogram shape identifies the failure
mode: a U-shape indicates an under-dispersed posterior (truth falls in
the tails too often), a central hump over-dispersion, and a monotone
slope a systematic bias.

\subsection{Joint coverage (TARP)}
\label{sec:validation:tarp}

Per-pixel ranks are insensitive to miscalibration of the
\emph{correlations} between dead pixels.  Joint calibration over the
$d = b^2$ box pixels is tested with the expected-coverage-probability
estimator of Ref.~\cite{lemos2023tarp} (TARP).  Coordinates are
standardized in log space using the pooled per-dimension mean and
standard deviation of all posterior samples.  For each image a random
reference point $\theta^{(i)}_r$ is drawn uniformly over the pooled
per-dimension range of samples and truths, and
\begin{equation}
  f_i \;=\; \frac{1}{J} \sum_{j=1}^{J}
    \mathbf{1}\!\left[\,
      \bigl\lVert x^{(i,j)} - \theta^{(i)}_r \bigr\rVert_2
      \;<\;
      \bigl\lVert x^{\mathrm{true}}_{i} - \theta^{(i)}_r \bigr\rVert_2
    \right].
  \label{eq:tarp-f}
\end{equation}
Under exact calibration $f_i \sim \mathcal{U}(0,1)$, equivalently the
expected coverage $\mathrm{ECP}(\alpha) = \Pr(f_i \le \alpha)$ lies on
the diagonal.  We report the ECP curve, its maximum absolute deviation
from the diagonal, and a Kolmogorov--Smirnov $p$-value for
$f \sim \mathcal{U}(0,1)$.  Under-dispersion appears as an ECP curve
sagging below the diagonal.

\subsection{Pulls and central credible intervals}
\label{sec:validation:pulls}

Two complementary, more interpretable diagnostics are evaluated in log
space.  The per-pixel pull
\begin{equation}
  z_{ip} = \frac{x^{\mathrm{true}}_{ip} - \widehat{\mu}_{ip}}
                {\widehat{\sigma}_{ip}},
  \qquad
  \widehat{\mu}_{ip} = \tfrac{1}{J}\textstyle\sum_j x^{(i,j)}_p,
  \quad
  \widehat{\sigma}_{ip}^2 =
      \tfrac{1}{J-1}\textstyle\sum_j
      \bigl(x^{(i,j)}_p - \widehat{\mu}_{ip}\bigr)^2,
  \label{eq:pull}
\end{equation}
is approximately standard normal for a calibrated, locally
near-Gaussian posterior; the pooled pull mean measures bias and the
pooled pull width the dispersion ratio (width $> 1$ $\Leftrightarrow$
under-dispersed posterior).  Independently, central credible intervals
at nominal levels $\gamma = 0.1, \dots, 0.9$ are formed from the
empirical $\bigl(\tfrac{1\mp\gamma}{2}\bigr)$-quantiles of the $J$
samples, and empirical coverage is compared to $\gamma$ both pooled over
dead pixels and for the physics observable
$\Sigma E = \sum_{p \in \mathrm{box}} E_p$ (in GeV), the latter probing
a physically relevant joint functional of the posterior.

\subsection{Comparison protocol and robustness}
\label{sec:validation:protocol}

All algorithms are evaluated on identical truth images, masks, and
per-image random streams, so cross-algorithm differences are paired.
The comparison figure overlays, per algorithm, the central-coverage
curves and the deviation of the empirical rank CDF from uniformity.
Statistical power at the study scale: pooled rank tests use
$N b^2$ ranks across $J{+}1$ bins; the TARP ECP has standard error
$\sim \sqrt{\alpha(1-\alpha)/N}$, so $N = 1000$ resolves coverage
distortions at the few-percent level.  Analysis-level robustness
safeguards: images containing non-finite samples are excluded and their
count reported (a nonzero count indicates a sampler failure, not an
analysis artifact), and partially completed runs are analyzed up to
their checkpoint, enabling monitoring of running studies.

\subsection{Generation-stage verification}
\label{sec:validation:generation}

Upstream of the inpainting study, the event-generation stage is verified
against Ref.~\cite{torbunov2024calo}: tower-area $E_T$ spectra
($1{\times}1$, $4{\times}4$, $7{\times}7$, $11{\times}11$, sliding
windows, periodic in $\phi$) and the $\Sigma E_T$ distribution per
centrality, overlaid across the five independently trained model seeds
with seed-ratio panels; the event-averaged
$\langle \sigma_{E_T}\rangle$ versus $\Sigma E_T$ profiles; regression
anchors on $\langle \Sigma E_T \rangle$ per centrality (soft $\pm 5\%$
checks against validated reference runs); and a monitor of the
high-energy single-tower tail reporting the rate of towers above
10~GeV per seed (towers of $\mathcal{O}(10\text{--}20)$~GeV are
kinematically plausible at $\sqrt{s_{NN}} = 200$~GeV; only energies
beyond 50~GeV are treated as unphysical).
